\section{The Federal Statutory Pathway}
\label{sec:federal}

\subsection{Structure of a Federal Mandate}

A federal statute mandating algorithmic redistricting would occupy the Elections Clause's second sentence: Congress exercises its power to ``make or alter'' state regulations of the manner of congressional elections. The operational structure requires minimal statutory text. The model statute in Section~\ref{sec:model} demonstrates that the core mandate can be stated in three operative provisions: a definition of the algorithmic method, a mandate that states apply that method to produce their congressional districts, and an enforcement provision conferring jurisdiction on federal district courts.

The statute need not specify every technical parameter of the algorithm---just as 2 U.S.C. \S\,2a does not derive the Huntington-Hill formula from first principles but instead codifies its essential mathematical operation. The statute should specify: (1) that the method must use only population data from the decennial census and geographic data from official federal sources; (2) that it must not use any data on voters' partisan affiliation, voting history, race, or ethnicity; (3) that it must optimize for population equality (within $\pm 0.5\%$) and geographic compactness; and (4) that its source code must be publicly available. Implementation details---the specific algorithm, software libraries, and parameter settings---can be delegated to administrative rulemaking by an appropriate federal agency, analogous to how the Census Bureau implements technical aspects of the apportionment computation.

\subsection{The Census Bureau's Role}

The Census Bureau (within the Department of Commerce) is the natural federal agency to administer algorithmic redistricting. The Bureau already produces the census data on which redistricting is based, maintains the TIGER geographic database from which district adjacency graphs are built, and has established relationships with state redistricting authorities. Its apportionment computation role provides a direct precedent for executing mathematical procedures that translate census data into politically consequential outputs.

Under the model statute, the Bureau would: (1) run the certified algorithm on final census data within 30 days of the apportionment announcement; (2) produce preliminary district maps for all 50 states; (3) publish those maps and the underlying code and data on a publicly accessible website; and (4) transmit the maps to state redistricting authorities as the baseline from which any minor adjustments may be made. States would have 60 days to identify specific communities of interest that require adjustment, subject to the constraints that adjustments must maintain population balance within the statutory tolerance and cannot reduce compactness below a specified threshold.

This structure preserves a meaningful role for state actors---they can propose adjustments for localized concerns---while ensuring that partisan data cannot be introduced through the adjustment process. The Bureau would promulgate rules specifying that proposed adjustments must be documented by reference to objective geographic or community factors, and that any adjustment that would predictably affect the partisan balance of a district is not a permissible basis.

\subsection{Comparison to the Huntington-Hill Statute}

The Huntington-Hill apportionment statute provides not merely an analogy but a drafting template. 2 U.S.C. \S\,2a(a) reads in full: ``The Bureau of the Census shall, in the same manner as provided in subsection (b), determine the number of persons in each State as shown by the decennial census of the population, and shall transmit to the Congress a statement showing the whole number of persons in each State, excluding Indians not taxed, as ascertained under the seventeenth and subsequent decennial censuses of the population, and the number of Representatives to which each State would be entitled under an apportionment of the then existing number of Representatives by the method known as the method of equal proportions, no State to receive less than one Member.''\footnote{2 U.S.C. \S\,2a(a).}

This statutory text mandates a specific mathematical method (equal proportions / Huntington-Hill), assigns implementation responsibility to the Census Bureau, specifies the data source (decennial census), and sets a minimum floor (one Representative per state). The algorithmic redistricting statute follows the same template: mandate a specific algorithmic method (recursive graph bisection via METIS or equivalent), assign implementation responsibility to the Census Bureau, specify the data sources (decennial census PL 94-171, TIGER shapefiles), and set constraints (population deviation $\leq \pm 0.5\%$, contiguity required).

The political history of apportionment illustrates why a statutory mathematical mandate is the right tool. Prior to 1941, Congress changed apportionment formulas in ways that benefited different states in different cycles, creating perennial controversy. After 1941, the Huntington-Hill method operated without congressional controversy for more than 80 years. The method's legitimacy rests not on its mathematical superiority over alternatives---scholars continue to debate which method is ``fairest'' in various senses---but on its procedural qualities: transparency, consistency, reproducibility, and resistance to manipulation.

\subsection{Dispute Resolution}

The model statute provides for expedited federal judicial review of challenges to algorithmic redistricting outputs. This differs from the pre-\textit{Rucho} partisan gerrymandering context in a critical way: the legal standard is administrable. A court reviewing whether an algorithmic redistricting output violates the statute need not engage in the unmanageable inquiry that deterred the \textit{Rucho} majority---it need not define what makes a partisan gerrymander excessive. The court need only determine whether the algorithm was applied as specified, whether the state's adjustment requests were evaluated under the statutory criteria, and whether the resulting districts satisfy population-equality requirements. These are straightforward factual and legal questions of the kind courts resolve routinely.

The model statute provides a three-judge district court panel for redistricting challenges, consistent with the tradition in redistricting litigation, with direct appeal to the Supreme Court. It imposes a 60-day deadline on the panel's decision to prevent drawn-out litigation from leaving states without valid districts for the next election cycle.
